%************************************************************************
% WuZ-spezifische Umgebungen fuer Skript, Folien und Handouts.
% Wird nach mathmacros(-sl) und macros geladen. Die Root-Dateien setzen
% \slidestrue (Folien) bzw. \handoutstrue (Handouts/Uebungsblaetter).
%
% Sichtbarkeit:
%   staunbox, merkbox            Skript, Folien, Handouts
%   vorbereitung, ablauf,
%   lehrhinweis                  nur Skript (Regie fuer Lehrende)
%   handout, uebungsblatt        Skript und Handouts (nicht Folien)
%   \nsl{...} / \os{...}         nur Skript / nur Folien (wie Embedded Computing)
%************************************************************************

% ---- Staunmoment: die Staun-Intervention der Lehreinheit -------------
\newtcolorbox{staunbox}[1]{enhanced,breakable,
  colback=rwuviolet!6,colframe=rwuviolet,boxrule=0.6pt,arc=2pt,
  fonttitle=\bfseries,coltitle=white,colbacktitle=rwuviolet,
  title={$\bigstar$\enspace Staunmoment: #1}}

% ---- Merksatz ---------------------------------------------------------
\newtcolorbox{merkbox}[1][Merke]{enhanced,breakable,
  colback=rwucyan40!25,colframe=rwucyan,boxrule=0.6pt,arc=2pt,
  fonttitle=\bfseries,coltitle=black,colbacktitle=rwucyan40,
  title={#1}}

% ---- Nur Skript: Vorbereitung, Ablauf, Hinweise fuer Lehrende ---------
\ifslides
  \excludeversion{vorbereitung}
  \excludeversion{ablauf}
  \excludeversion{lehrhinweis}
\else
  \newtcolorbox{vorbereitung}[1][]{enhanced,breakable,
    colback=black!3,colframe=rwugray,boxrule=0.5pt,arc=2pt,
    borderline west={3pt}{0pt}{rwucyan},
    fonttitle=\bfseries,coltitle=white,colbacktitle=rwugray,
    title={Vorbereitung vor der Stunde\ifx&#1&\else: #1\fi}}
  \newtcolorbox{ablauf}[1][Ablauf der Doppelstunde (90 Min.)]{enhanced,breakable,
    colback=rwucyan40!12,colframe=rwucyandark,boxrule=0.5pt,arc=2pt,
    fonttitle=\bfseries,coltitle=white,colbacktitle=rwucyandark,
    title={#1}}
  \newtcolorbox{lehrhinweis}[1][Hinweis für Lehrende]{enhanced,breakable,
    colback=rwuviolet!3,colframe=rwuvioletlight,boxrule=0.5pt,arc=2pt,
    fonttitle=\bfseries\small,coltitle=black,colbacktitle=rwuvioletlight!40,
    title={#1}}
\fi
\ifhandouts
  % Handouts gehen an Studierende: keine Regie, keine Loesungshinweise
  \excludeversion{lehrhinweis}
\fi

% ---- Checkliste in Vorbereitungsboxen -----------------------------------
\newenvironment{checkliste}{\begin{list}{$\square$}{\leftmargin 1.4em \labelwidth 1em
  \labelsep 0.4em \itemsep 0.1em \parsep 0em \topsep 0.3em}}{\end{list}}

% ---- Handouts und Uebungsblaetter ---------------------------------------
% Jede Lehreinheit hat eine Datei handouts/LEnn.tex mit einem Handout und
% einem Uebungsblatt in zwei Teilen. Im Skript wird sie am Kapitelende
% eingebunden, in handouts.tex als Druckvorlage, auf den Folien gar nicht.
%   uebungsblatt{Titel}           Teil A: wird zu Beginn der Stunde ausgeteilt,
%                                 darf den Staunmoment nicht vorwegnehmen
%   uebungsblattb{n}{Titel}       Teil B: nach dem Staunmoment; n = Nummer der
%                                 ersten Aufgabe (zaehlt Teil A weiter, auch
%                                 wenn Teil B allein gedruckt wird)
%   uebungsblattc{Titel}          Teil C (optional): Material zum Austeilen,
%                                 etwa die Quellen fuer eine Gruppenaufgabe
% Beschriftungen; handouts.tex ueberschreibt sie fuer die englische Fassung
\providecommand{\wuzkurs}{Wissenschaftliches Arbeiten (WuZ)}
\providecommand{\wuzlename}{Lehreinheit}
\providecommand{\wuzblattname}{Übungsblatt}
\providecommand{\wuzaufgabename}{Aufgabe}
\providecommand{\wuzpersonname}{Name}
\providecommand{\wuzdatumname}{Datum}
\newcounter{wuzaufgabe}
\newcommand{\wuzblattkopf}[2]{%
  \clearpage
  \noindent{\sffamily
  \colorbox{#1}{\parbox{\dimexpr\textwidth-2\fboxsep\relax}{\color{white}\small\bfseries
    \wuzkurs\hfill\wuzlename~\thechapter}}\\[3mm]
  {\LARGE\bfseries\color{#1}#2}\\[-1mm]
  {\color{#1}\rule{\textwidth}{0.5mm}}}\par\medskip}
\newenvironment{handout}[1]{\wuzblattkopf{rwuviolet}{Handout: #1}}{\par}
\newenvironment{uebungsblatt}[1]{\setcounter{wuzaufgabe}{0}%
  \wuzblattkopf{rwucyandark}{\wuzblattname~A: #1}%
  \noindent{\small \wuzpersonname: \rule{55mm}{0.3pt}\hfill \wuzdatumname: \rule{30mm}{0.3pt}}\par\medskip}{\par}
\newenvironment{uebungsblattb}[2]{\setcounter{wuzaufgabe}{\numexpr#1-1\relax}%
  \wuzblattkopf{rwucyandark}{\wuzblattname~B: #2}%
  \noindent{\small \wuzpersonname: \rule{55mm}{0.3pt}\hfill \wuzdatumname: \rule{30mm}{0.3pt}}\par\medskip}{\par}
\newenvironment{uebungsblattc}[1]{\wuzblattkopf{rwucyandark}{\wuzblattname~C: #1}}{\par}
% Volltexte zu Blatt C (LE 4); im Skript leer, handouts.tex definiert sie
\providecommand{\LEvierVolltexte}{}
% \aufgabe{Titel}[Bearbeitungszeit/Form]
\NewDocumentCommand{\aufgabe}{m o}{\par\medskip\refstepcounter{wuzaufgabe}%
  \noindent{\bfseries\color{rwucyandark}\wuzaufgabename~\thewuzaufgabe: #1}%
  \IfValueT{#2}{\hfill{\small\color{rwugray}#2}}\par\smallskip}
% Schreiblinien fuer Antworten: \schreiblinien{Anzahl}
\newcommand{\schreiblinien}[1]{\par\vspace{2mm}\foreach \i in {1,...,#1}{%
  \noindent\color{rwugray!60}\rule{\textwidth}{0.3pt}\par\vspace{6mm}}\color{black}}

\ifslides
  \newcommand{\handoutinput}[1]{}
\else
  \newcommand{\handoutinput}[1]{\input{#1}}
\fi
% in handouts.tex: \handoutle{Nummer}{Datei}
\newcommand{\handoutle}[2]{\setcounter{chapter}{#1}\input{#2}}

% ---- Kleinkram -------------------------------------------------------------
\newcommand{\LE}[1]{LE~#1}
\newcommand{\buch}[1]{\emph{WuZ-Buch}, #1}   % Verweis auf das Begleitbuch
